\chapter{Замыкание Hodge-родителя с одной размерной калибровкой}
\label{ch:v7-single-scale-calibration-closure-gate}

\section{Честный EFT-вопрос}

Предыдущий гейт запретил выводить $\mu$ из одного фона $H_{15}$. Теперь
принимается минимальная разрешённая альтернатива: одна физическая масса
используется как единственный размерный вход. Нужно проверить, становятся ли
после этого все остальные величины предсказаниями.

На пространстве одиннадцати комплексных рёбер рассмотрим лоренцево
допустимую эффективную плотность
\begin{equation}
 \mathcal L_E
 =\frac{Z}{2}\sum_e\left[(\partial_\mu x_e)^2
 +(\partial_\mu y_e)^2\right]
 -\kappa\mathcal S_\mu(z),
 \qquad z_e=x_e+iy_e,
 \label{eq:v7-single-scale-effective-lagrangian}
\end{equation}
где $Z,\kappa>0$ общи для всех рёбер, а $\mathcal S_\mu$ является уже
выведенным Hodge-потенциалом. Отдельные $Z_e$, $\kappa_e$ и $\mu_e$
запрещены.

\section{Канонические переменные}

После замены
\begin{equation}
 \phi_e=\sqrt Z\,z_e,
 \qquad
 v=\sqrt Z\,\mu,
 \qquad
 \lambda_E=\frac{\kappa}{Z^2}
 \label{eq:v7-single-scale-canonical-variables}
\end{equation}
весь квадратичный спектр зависит от одной размерной комбинации
\begin{equation}
 M_0^2=\frac{\kappa\mu^2}{Z}=\lambda_Ev^2.
 \label{eq:v7-single-scale-common-mass}
\end{equation}
Это показывает, что отдельные $\mu$, $Z$ и $\kappa$ не являются тремя
размерными входами. Линейная физика видит только $M_0$.

\section{Безразмерное предсказание спектра}

В нуле обобщённый гессиан в единицах $M_0^2$ равен
\begin{equation}
 \Spec(M_0^{-2}G^{-1}\Hess_0)
 =\{-4^{(12)},\;4^{(10)}\}.
 \label{eq:v7-single-scale-origin-spectrum}
\end{equation}
В Hodge-вакууме получаем
\begin{equation}
 \Spec(M_0^{-2}G^{-1}\Hess_*)
 =\{0^{(6)},\;4^{(10)},\;8^{(6)}\}.
 \label{eq:v7-single-scale-vacuum-spectrum}
\end{equation}
Шесть нулей являются фазами выбранных рёбер. Десять компонент пяти
нежелательных рёбер имеют массу $2M_0$, а шесть выбранных радиальных мод ---
$2\sqrt2M_0$. Следовательно,
\begin{equation}
 \frac{m_{\rm rad}}{m_{\rm gap}}=\sqrt2,
 \qquad
 \frac{\xi_{\rm rad}}{\xi_{\rm gap}}=\frac1{\sqrt2}.
 \label{eq:v7-single-scale-linear-ratios}
\end{equation}
Одна измеренная ненулевая масса действительно калибрует весь внутренний
линейный спектр. Отношение $\sqrt2$ не использует абсолютную единицу.

\begin{theorem}[Частичное одномасштабное замыкание]
При единой положительной кинетической метрике одна размерная калибровка
фиксирует все ненулевые собственные значения линейного Hodge-спектра и его
корреляционные длины. В частности, отношение выбранной радиальной массы к
щели нежелательного ребра равно $\sqrt2$.
\end{theorem}

\section{Почему полный EFT ещё не замкнут}

Одна масса фиксирует $M_0$, но не разлагает его на $v$ и $\lambda_E$:
\begin{equation}
 v=\frac{M_0}{\sqrt{\lambda_E}}.
 \label{eq:v7-single-scale-vev-coupling-degeneracy}
\end{equation}
При фиксированном $M_0$ значения
\begin{equation}
 \lambda_E\in\left\{\frac14,1,4\right\}
 \quad\Longrightarrow\quad
 v\in\left\{2M_0,M_0,\frac12M_0\right\}
 \label{eq:v7-single-scale-degeneracy-example}
\end{equation}
дают один и тот же квадратичный спектр, но разные четырёхточечные вершины и
нелинейные энергии. Например, натяжение профиля коразмерности один имеет
масштаб
\begin{equation}
 \sigma_{\rm nl}\sim v^2M_0=\frac{M_0^3}{\lambda_E},
 \label{eq:v7-single-scale-nonlinear-energy}
\end{equation}
которая не определяется одной массовой калибровкой.

Конечный полевой след ранее дал общую положительную метрику $3I$, но
коэффициент её пространственно-временного теплового продолжения не выведен.
Это согласуется со спектральным действием: функция отсечения и масштаб входят
в исходный функционал, а связи кинетического и скалярного секторов требуют
конкретного произведённого спектрального оператора
\cite{v7-single-scale-spectral-action,v7-single-scale-product}.

\begin{nogocriterion}[Запрет полного одномасштабного статуса]
Нельзя называть отношение $\sqrt2$ полным физическим замыканием, пока
$\lambda_E=\kappa/Z^2$ не выведено из одного пространственно-временного
спектрального следа. Выбор $Z=1$ или $\lambda_E=1$ является выбором
нормировки, а не вторым предсказанием. Массы линейных мод замкнуты условно;
амплитуды, рассеяние и нелинейные энергии остаются открыты.
\end{nogocriterion}

\begin{protocol}
\begin{enumerate}
 \item размерных калибровочных входов: \textbf{$1$};
 \item общий квадратичный масштаб: \textbf{$M_0^2=\kappa\mu^2/Z$};
 \item радиальных выбранных мод: \textbf{$6$ с массой $2\sqrt2M_0$};
 \item нежелательных массивных компонент: \textbf{$10$ с массой $2M_0$};
 \item фазовых нулей: \textbf{$6$};
 \item предсказанное отношение масс: \textbf{$\sqrt2$};
 \item линейные корреляционные длины: \textbf{замкнуты одной массой};
 \item эффективная квартика $\lambda_E$: \textbf{не выведена};
 \item нелинейные энергии и амплитуды: \textbf{не замкнуты};
 \item физическая временная нормировка: \textbf{не выведена};
 \item итог: \textbf{частичное линейное одномасштабное замыкание}.
\end{enumerate}
\end{protocol}

Машинный сертификат сохранён в
\path{s2t_v7_single_scale_calibration_closure_gate_results.json}.

\begin{thebibliography}{9}
\bibitem{v7-single-scale-spectral-action}
A.~H.~Chamseddine, A.~Connes,
\emph{The Spectral Action Principle},
Communications in Mathematical Physics 186 (1997), 731--750;
arXiv:hep-th/9606001.
\bibitem{v7-single-scale-product}
S.~Brain, B.~Mesland, W.~D.~van Suijlekom,
\emph{Gauge Theory for Spectral Triples and the Unbounded Kasparov Product},
Journal of Noncommutative Geometry 10 (2016), 135--206;
arXiv:1306.1951.
\end{thebibliography}